\documentclass[11pt]{article}
\usepackage[T1]{fontenc}
\usepackage{lmodern}
\usepackage[margin=1in]{geometry}
\usepackage{amsmath,amssymb,amsthm}
\usepackage{microtype}
\usepackage[colorlinks=true,linkcolor=blue,citecolor=blue,urlcolor=blue]{hyperref}
\newtheorem{theorem}{Theorem}
\newtheorem{lemma}[theorem]{Lemma}
\newtheorem{proposition}[theorem]{Proposition}
\newcommand{\C}{\mathbb C}
\newcommand{\R}{\mathbb R}
\newcommand{\B}{\mathcal B}
\newcommand{\dz}{\partial_z}
\newcommand{\db}{\partial_{\bar z}}
\newcommand{\pv}{\operatorname{pv}}
\DeclareMathOperator{\supp}{supp}
\title{A note on the sharp Beurling--Ahlfors bound}
\author{Kaiwen Jin\\[2pt]\small\href{mailto:yui.ui.siki@gmail.com}{yui.ui.siki@gmail.com}}
\date{October 12, 2026}
\begin{document}
\maketitle

\begin{abstract}
Agazzi and collaborators prove the sharp Beurling--Ahlfors norm in
arXiv:2609.28311v2. We give an explicit derivation of the corresponding
disk principal-value estimate for arbitrary smooth compactly supported
complex functions. The argument starts from their unrestricted
quasiconvexity theorem and records the power integral, the Cauchy-potential
cutoff, and complex duality. This is an attributed exposition of the
existing result.
\end{abstract}

\section{The estimate and its source}

Identify \(\C\) with \(\R^2\), write \(dA=dx\,dy\), and use
\[
 \dz=\tfrac12(\partial_x-i\partial_y),\qquad
 \db=\tfrac12(\partial_x+i\partial_y).
\]
For a smooth compactly supported complex function \(f\), the
Beurling--Ahlfors transform is
\begin{equation}\label{eq:pv}
 \B f(z)=-\frac1\pi\lim_{\epsilon\downarrow0}
       \int_{|z-w|>\epsilon}\frac{f(w)}{(z-w)^2}\,dA(w).
\end{equation}
All \(L^p\) norms below use complex modulus and area measure.

\begin{theorem}\label{thm:bound}
For every real \(1<p<\infty\) and every \(f\in C_c^\infty(\C;\C)\),
the limit in \eqref{eq:pv} exists at every point, \(\B f\in L^p(\C)\),
and
\begin{equation}\label{eq:bound}
 \|\B f\|_p\le (p^*-1)\|f\|_p,\qquad
 p^*=\max\{p,p/(p-1)\}.
\end{equation}
\end{theorem}

The result is due to Agazzi et al.~\cite{agazzi}: their Corollary~2,
equation~(4), states the full complex \(L^p\) operator norm equality.
The purpose of this note is to spell out the passage to
\eqref{eq:pv} and \eqref{eq:bound}. In particular, neither component of
\(f\) is restricted, and its integral need not vanish.

We use the following precise form of their main theorem. Write a
real-linear map \(A:\C\to\C\) as
\[
 Az=a_+z+a_-\bar z,\qquad a=|a_+|,\quad b=|a_-|,\quad s=a+b.
\]
Then \(|A|_{\mathrm{op}}=s\) and \(\det A=a^2-b^2\). Define
\begin{equation}\label{eq:L}
 L(A)=
 \begin{cases}
   \det A,&s\le1,\\
   2a-1,&s>1 .
 \end{cases}
\end{equation}
The two expressions agree when \(s=1\).

\begin{theorem}[Agazzi et al., Theorem 1]\label{thm:source}
On \(Q=\R^2/\mathbb Z^2\), with area one, every
\(A\in\R^{2\times2}\) and every
\(\varphi\in C^\infty_{\mathrm{per}}(Q;\R^2)\) satisfy
\begin{equation}\label{eq:qc}
 \int_Q L(A+D\varphi)\,dx\ge L(A).
\end{equation}
\end{theorem}

Theorem~\ref{thm:source} is quoted from \cite[Theorem~1 and
Section~2]{agazzi}. It concerns arbitrary Jacobians of two-component
fields. The source proves it by a constrained gradient projection and a
pointwise representation of \(L\). We need only its stated inequality.
The relation between \(L\), power integrands, and the transform was
already considered by Baernstein and Montgomery-Smith~\cite{bms};
we include the calculation needed here.

\section{From periodic quasiconvexity to power estimates}

Let \(u\in C_c^\infty(\R^2;\R^2)\). Choose a square
\(x_0+[0,R]^2\) containing its support in the interior. Then
\(\varphi(y)=u(x_0+Ry)/R\) extends smoothly and periodically to \(Q\),
and \(D\varphi(y)=Du(x_0+Ry)\). Theorem~\ref{thm:source} at \(A=0\)
gives
\[
 \int_{\R^2}L(Du)\,dx=R^2\int_Q L(D\varphi)\,dy\ge0.
\]
Also \(\int_{\R^2}\det Du=0\): both
\(\int u^1_xu^2_y\) and \(\int u^1_yu^2_x\) equal
\(-\int u^1u^2_{xy}\) by integration by parts. Thus, with \(M=L-\det\),
\begin{equation}\label{eq:Mpositive}
 \int_{\R^2}M(Du/t)\,dx\ge0\qquad(t>0),
\end{equation}
where we have applied the same argument to \(u/t\).

Fix \(p>2\). For a matrix \(A\) with \(s=a+b>0\), \(M(A/t)\)
vanishes when \(t\ge s\), whereas for \(0<t<s\) it is
\((b^2-a^2)t^{-2}+2at^{-1}-1\). Consequently
\begin{align}
 \int_0^\infty t^{p-1}M(A/t)\,dt
 &=\frac{(b^2-a^2)s^{p-2}}{p-2}
      +\frac{2as^{p-1}}{p-1}-\frac{s^p}{p}\notag\\
 &=\frac{2}{p(p-1)(p-2)}
                [(p-1)b-a](a+b)^{p-1}. \label{eq:mellin}
\end{align}
The last equality uses \(b^2-a^2=(b-a)s\).
Both sides vanish at \(A=0\). Moreover,
\begin{equation}\label{eq:absolute}
 \int_0^\infty t^{p-1}|M(A/t)|\,dt
 \le\left(\frac1{p-2}+\frac2{p-1}+\frac1p\right)s^p .
\end{equation}
Indeed, bound \(|b^2-a^2|\) by \(s^2\) and \(a\) by \(s\)
in the three integrals. Since \(Du\) is bounded and compactly supported,
\eqref{eq:absolute} justifies Fubini. Integrating
\eqref{eq:Mpositive} with the positive weight \(t^{p-1}\) therefore gives
\begin{equation}\label{eq:Burkholder}
 \int_{\R^2}[(p-1)b-a](a+b)^{p-1}\,dx\ge0,
 \qquad a=|\dz u|,\quad b=|\db u|.
\end{equation}
Here \(u=u^1+iu^2\); the complex coefficients of its real Jacobian
are exactly \(\dz u\) and \(\db u\).

\begin{lemma}\label{lem:majorant}
If \(p>2\), \(k=p-1\), and
\(\alpha=k^{p-1}/p^{p-2}\), then for every \(a,b\ge0\),
\begin{equation}\label{eq:majorant}
 \alpha(kb-a)(a+b)^{p-1}\le k^p b^p-a^p .
\end{equation}
\end{lemma}

\begin{proof}
For \(b>0\), divide by \(b^p\) and put \(t=a/b\).
The difference between the right and left sides becomes
\[
 E(t)=k^p-t^p-\alpha(k-t)(t+1)^{p-1}.
\]
We have \(E(k)=0\), and
\[
 E'(t)=p\bigl[-t^{p-1}
          +\alpha(t+1)^{p-2}(t-(p-2))\bigr].
\]
The derivative is negative for \(0\le t\le p-2\).
For \(t>p-2\), set
\[
 R(t)=\frac{t^{p-1}}{(t+1)^{p-2}(t-(p-2))}.
\]
Then \(R(k)=\alpha\), and
\[
 \frac{R'(t)}{R(t)}
 =\frac{p-1}{t}-\frac{p-2}{t+1}-\frac1{t-(p-2)}
 =-\frac{(p-1)(p-2)}{t(t+1)(t-(p-2))}<0.
\]
Thus \(E\) decreases up to \(k\) and increases thereafter.
Its minimum is \(E(k)=0\). Continuity as \(b\downarrow0\)
proves the remaining case.
\end{proof}

Integrating \eqref{eq:majorant}, then using \eqref{eq:Burkholder},
proves the following potential estimate:
\begin{proposition}\label{prop:potential}
For every \(p>2\) and every \(u\in C_c^\infty(\C;\C)\),
\begin{equation}\label{eq:potential}
 \|\dz u\|_p\le(p-1)\|\db u\|_p .
\end{equation}
\end{proposition}
There is no restriction on either real component of \(u\).

\section{The Cauchy potential and the disk kernel}

The locally integrable function \(E_C(z)=1/(\pi z)\) satisfies
\begin{equation}\label{eq:fundamental}
 \db E_C=\delta_0,\qquad
 \dz E_C=-\pv\frac1{\pi z^2},
\end{equation}
where the principal value uses disks. We record the boundary computation
to fix the sign and normalization. For
\(\psi\in C_c^\infty(\C;\C)\), polar coordinates give
\[
 \db\psi=\tfrac12e^{i\theta}(\psi_r+i\psi_\theta/r),
 \qquad
 \dz\psi=\tfrac12e^{-i\theta}(\psi_r-i\psi_\theta/r).
\]
Integration in \(r>\epsilon\), and periodic integration by parts in
\(\theta\), yield
\begin{align*}
 -\int_{|z|>\epsilon} E_C\,\db\psi\,dA
   &=\frac1{2\pi}\int_0^{2\pi}\psi(\epsilon e^{i\theta})\,d\theta,\\
 -\int_{|z|>\epsilon} E_C\,\dz\psi\,dA
   &=\frac1{2\pi}\int_0^{2\pi}e^{-2i\theta}
                         \psi(\epsilon e^{i\theta})\,d\theta
       -\int_{|z|>\epsilon}\frac{\psi(z)}{\pi z^2}\,dA.
\end{align*}
The first boundary expression tends to \(\psi(0)\).
In the second, the constant term has zero angular integral, and the
remainder is \(O_\psi(\epsilon)\). Subtracting \(\psi(0)\) near zero
also proves existence of the principal value, since
\(|\psi(z)-\psi(0)|\le C_\psi|z|\). The omitted disk integrals on the
left tend to zero because \(1/|z|\) is locally integrable.
This proves \eqref{eq:fundamental}, without an additional delta term.

For \(f\in C_c^\infty(\C;\C)\), put
\[
 u(z)=(E_C*f)(z)=\int_\C\frac{f(w)}{\pi(z-w)}\,dA(w).
\]
It is smooth: after writing the integral as
\(\int E_C(w)f(z-w)\,dA(w)\), place each derivative on \(f\).
On every compact \(z\)-set, the resulting integrals have common bounded
support in \(w\) and an integrable bound proportional to \(|E_C(w)|\).
Dominated convergence applies to every derivative.
Convolving \eqref{eq:fundamental} with \(f\) gives
\begin{equation}\label{eq:potentialidentity}
 \db u=f,\qquad \dz u=\B f .
\end{equation}
The second equality holds for the pointwise definition
\eqref{eq:pv}: subtract \(f(z)\) in a disk about \(z\); the constant
has zero angular integral, and the error from a removed disk of radius
\(\epsilon\) is \(O_f(\epsilon)\), locally uniformly in \(z\).
Thus the principal values converge to the smooth distributional derivative.
This is also the classical normalization in
\cite[equations~(1.3), (2.1), (2.2)]{bms}.

Choose \(R_0>0\) with \(\supp f\subset B_{R_0}\).
For \(|z|>2R_0\), differentiating the nonsingular integral gives
\[
 |u(z)|\le C_f|z|^{-1},\qquad |Du(z)|\le C_f|z|^{-2}.
\]
Local smoothness and this decay imply \(Du\in L^r\) for every \(r>1\).
The potential itself need not belong to \(L^2\).

Let \(0\le\chi\le1\) be smooth, equal to one on \(B_1\), zero outside
\(B_2\), and put \(u_R(z)=\chi(z/R)u(z)\).
For every fixed \(r>1\),
\begin{equation}\label{eq:cutoff}
 Du_R-Du=(\chi(z/R)-1)Du+R^{-1}u(z)D\chi(z/R)
                   \longrightarrow0\quad\hbox{in }L^r .
\end{equation}
The first term tends to zero by integrability of \(|Du|^r\).
The second is supported in \(R\le|z|\le2R\), is bounded by
\(C_fR^{-2}\), and has norm at most \(C_{f,r}R^{2/r-2}\),
which tends to zero. Applying Proposition~\ref{prop:potential} to
\(u_R\) and passing in norm proves \eqref{eq:bound} for every \(p>2\).

\section{The remaining exponents}

Use \(\widehat h(\xi)=\int_{\R^2}h(x)e^{-2\pi i x\cdot\xi}\,dx\)
and its unitary \(L^2\) extension. For the compact smooth \(u_R\),
integration by parts gives
\[
 \widehat{\dz u_R}=\pi i(\xi_1-i\xi_2)\widehat{u_R},\qquad
 \widehat{\db u_R}=\pi i(\xi_1+i\xi_2)\widehat{u_R}.
\]
Define
\begin{equation}\label{eq:multiplier}
 m(\xi)=\frac{\xi_1-i\xi_2}{\xi_1+i\xi_2}\quad(\xi\ne0),\qquad
 m(0)=1.
\end{equation}
It has unit modulus and is even. The preceding identities give
\(\widehat{\dz u_R}=m\widehat{\db u_R}\) almost everywhere.
By \eqref{eq:cutoff} at \(r=2\) and Plancherel,
\[
 \widehat{\B f}=m\widehat f,\qquad \|\B f\|_2=\|f\|_2.
\]
This identifies the multiplier with the original disk operator.

On complex \(L^2\), let \(B_2=\mathcal F^{-1}m\mathcal F\) and
use the pairing \(\langle h,g\rangle=\int h\bar g\,dA\).
Its adjoint has multiplier \(\bar m\).
For the antilinear isometry \(Ch=\bar h\),
\(\widehat{Ch}(\xi)=\overline{\widehat h(-\xi)}\);
evenness of \(m\) therefore gives \(B_2^*=CB_2C\).

Fix \(1<q<2\) and set \(p=q/(q-1)>2\). For every
\(g\in C_c^\infty(\C;\C)\), the bound already proved at exponent \(p\)
implies
\[
 \|B_2^*g\|_p=\|\B(Cg)\|_p\le(p-1)\|g\|_p .
\]
For the original smooth compact \(f\), the \(L^2\) adjoint identity
and Hölder give
\begin{equation}\label{eq:pairing}
 \left|\int\B f\,\bar g\,dA\right|
 =\left|\int f\,\overline{B_2^*g}\,dA\right|
 \le(p-1)\|f\|_q\|g\|_p .
\end{equation}
We already know \(\B f\in L^q\) from its smoothness and decay.
The complex \(L^q\) norm is the supremum of its pairings with unit
\(L^p\) test functions, and smooth compact tests suffice by density.
For example, for a nonzero \(h\in L^q\) the unit norm maximizer is
\(h|h|^{q-2}/\|h\|_q^{q-1}\), set to zero where \(h=0\).
Approximation in \(L^p\) and Hölder recover its pairing from smooth tests.
If \(g_n\) approximates the unit test, replacing it by
\(g_n/\max\{1,\|g_n\|_p\}\) preserves the limiting pairing.
For density, first truncate an \(L^p\) function to a bounded compactly
supported \(v\), with \(|v|\le M\). Theorem~1.3.20(iii) of \cite{tao}
provides a continuous compactly supported approximation in \(L^1\).
Radially projecting this approximant onto the complex disk of radius \(M\)
gives \(w\) with \(|w|\le M\), preserving continuity and compact support;
its \(L^1\) error increases by at most a factor of two. Hence
\[
 \|v-w\|_p^p\le(2M)^{p-1}\|v-w\|_1.
\]
Choose the \(L^1\) error after \(M\). Mollification of \(w\) converges
uniformly on a common compact support, hence in \(L^p\); truncation
converges in \(L^p\) by dominated convergence. This proves the required
smooth density without restricting the complex phase.
Taking the supremum in \eqref{eq:pairing} yields
\[
 \|\B f\|_q\le(p-1)\|f\|_q
            =\frac1{q-1}\|f\|_q .
\]
Together with the \(p>2\) estimate and \(p=2\) equality, this proves
Theorem~\ref{thm:bound} for every \(1<p<\infty\), with
\(\max\{p,p/(p-1)\}-1\) and the kernel \(-1/(\pi z^2)\).
All inputs throughout are arbitrary complex smooth compactly supported
functions on the plane.

\begin{thebibliography}{9}
\bibitem{agazzi}
A. Agazzi, K. Astala, R. Ba\~nuelos, G. Bruno, G. Cassese,
D. Faraco, A. Guerra, M. Hormozi, B. Kirchheim, A. Koski,
V. Kova\v c, J. Kristensen, F. Pasqualotto, I. Prause,
E. Saksman, V. \v Sver\'ak, L. Sz\'ekelyhidi, R. Tione,
A. Volberg, and H. Yang,
\emph{Sharp Higher Integrability Theory, Part I:
The \(L^p\)-norm of the Beurling--Ahlfors transform},
arXiv:2609.28311v2 (2026).
\href{https://arxiv.org/abs/2609.28311v2}{Versioned primary source};
\href{https://arxiv.org/html/2609.28311v2}{complete proof}.

\bibitem{bms}
A. Baernstein II and S. J. Montgomery-Smith,
\emph{Some conjectures about integral means of \(\partial f\) and
\(\bar\partial f\)},
in \emph{Complex Analysis and Differential Equations},
C. Kiselman (ed.), Acta Universitatis Upsaliensis C,
64 (1999), 92--109.
\href{https://stephenmontgomerysmith.github.io/preprints/brn.html}
{Author's publication page};
\href{https://stephenmontgomerysmith.github.io/preprints/brn6.tex}
{Author's TeX source}.

\bibitem{tao}
T. Tao, \emph{An Introduction to Measure Theory},
Graduate Studies in Mathematics 126, American Mathematical Society, 2011.
\href{https://terrytao.wordpress.com/wp-content/uploads/2012/12/gsm-126-tao5-measure-book.pdf}
{Author's preliminary version}.
\end{thebibliography}

\paragraph{AI-use disclosure.}
Codex assisted with primary-source comparison, the displayed transfer
calculations, drafting, and document checks. The sharp norm result and its
quasiconvexity proof are attributed to Agazzi and collaborators.
Independent mathematical review of this exposition is requested.

\end{document}
